\section{Discussion and Future Work}
\label{sec:discussion}

This paper establishes a pilot baseline with two outputs: a feasibility
profile for the generation-and-audit pipeline, and a set of preliminary
proxy-based directional patterns to prioritize for future human
evaluation. The first is a generator-feasibility result: roadway
markings and maintenance interventions realize reliably across cities,
lighting repair remains the clearest bottleneck (valid rate 0.30), and
the null baseline-score relationship (Spearman $\rho{=}0.10$) suggests
that multi-lever richness is driven more by scene geometry and visible
editable supports than by starting safety level. The second layer is
more substantive: Mobility
Infrastructure produces the largest retained auxiliary safety shifts
under the current Place Pulse proxy. Whether that reflects genuine human
perceptual sensitivity to roadway-legibility cues, scorer-specific
priors, or an interaction with edit realizability remains open.

The heterogeneity within Environmental Amenity is equally informative.
Localized greenery addition is strongly positive, while lighting repair
and canopy management are weak or negative on average despite being
theory-grounded cues. That pattern suggests family-level aggregation can
hide important lever-specific differences in editability and scorer
response, and motivates finer-grained prioritization at the lever level
rather than the family level alone.

\noindent\textbf{On the auxiliary proxy.} The auxiliary scorer
$f_a$ has not been validated on diffusion-edited images, and
$\Delta_{\mathrm{aux}}$ should not be read as a calibrated perceptual
effect size. The claim here is narrower: for semantically targeted,
monotone interventions such as repainting markings, removing surface
disorder, or adding greenery, the directional signal is grounded in a
model trained on large-scale crowdsourced pairwise human preference over
closely related scene variation, and the sign and relative magnitude of
$\Delta_{\mathrm{aux}}$ are useful for prioritizing which retained edits
to send to human evaluation first.

\noindent\textbf{XAI contribution.} Beyond urban perception, the main
contribution is methodological. The framework reframes visual
explanation around structured, theory-grounded interventions rather than
post hoc attribution over pixels or regions, and it imposes a stronger
standard in which a proposed change must be semantically coherent and
generatively feasible before entering the attribution result.

\noindent\textbf{Future directions and human endpoint.} The present
pipeline is intentionally minimal: prompt-only generation, a prompted
VLM critic, and an auxiliary proxy scorer establish a reproducible
lower-bound baseline rather than a fully validated explanation system.
The next steps follow directly from that structure: planner triage to
allocate the candidate budget more intelligently, stronger generation
control for validity, critic calibration against human judgements,
pairwise human evaluation as the ground-truth endpoint, and eventually
end-to-end optimization once planner, generator, and critic are each
validated independently. Concretely, the 9-page analysis points to four
near-term workstreams: predicting lever viability from scene features so
the candidate budget is spent more selectively; improving generator
alignment on known failure modes such as lighting repair; calibrating
the critic against human E1 raters and 2AFC judgements; and only then
considering joint optimization across planner, generator, and critic.
